% ============================================================================
%  SECCIÓN 1.3: PLANTEAMIENTO DEL PROBLEMA
% ============================================================================

\section{Planteamiento del problema}

\subsection{Situaci\'on problem\'atica}

A partir de lo expuesto en los antecedentes, la \textbf{situaci\'on
problem\'atica} se puede establecer de la siguiente manera: la informaci\'on
sobre los escenarios deportivos a nivel nacional se encuentra dispersa,
desactualizada y sin un medio unificado de consulta, lo que impide que la
poblaci\'on conozca, localice y acceda a la oferta deportiva y limita la
visibilidad de los propios recintos.

Las causas de esta situaci\'on son la ausencia de un registro centralizado de
escenarios deportivos, la escasa digitalizaci\'on del sector y la falta de
herramientas de visualizaci\'on interactiva adaptadas al contexto nacional.
Como efecto, se observa una baja asistencia a los eventos deportivos, la
subutilizaci\'on de la infraestructura existente y el desaprovechamiento de
las oportunidades econ\'omicas y sociales que el deporte genera. Como se\~nala
\citet{arias2012}, el problema constituye el punto de partida de toda
investigaci\'on y surge cuando se identifica una situaci\'on de carencia o una
necesidad no satisfecha, tal como ocurre en este caso.

\subsection{Objeto de estudio}

El objeto de estudio del presente proyecto es el \textbf{proceso de difusi\'on
y acceso a la informaci\'on sobre los escenarios deportivos a nivel
nacional}, incluyendo los mecanismos de registro, actualizaci\'on y consulta
de los datos de cada recinto. Se justifica su estudio porque la mejora de
este proceso, mediante una soluci\'on tecnol\'ogica, permite incrementar la
pr\'actica deportiva de la poblaci\'on, optimizar la gesti\'on de los espacios
deportivos y favorecer la toma de decisiones de las entidades
administradoras.

\subsection{Estudio de soluciones}

Se analizaron soluciones contempor\'aneas aplicadas en contextos similares,
tanto nacionales como internacionales, con el fin de identificar las
pr\'acticas m\'as efectivas. La siguiente tabla presenta la valoraci\'on de los
casos de estudio analizados:

\begin{table}[H]
  \centering
  \caption{Valoraci\'on de soluciones existentes}
  \label{tab:estudio-soluciones}
  \begin{tablaAPA}
    \begin{tabular}{@{}p{4.2cm}p{3.0cm}p{4.2cm}p{3.6cm}@{}}
      \toprule
      \textbf{Casos de estudio} & \textbf{Problema} &
      \textbf{Soluci\'on aplicada} & \textbf{Valoraci\'on} \\
      \midrule
      Aplicaciones internacionales de mapas deportivos &
      Informaci\'on deportiva dispersa &
      Sistema web y m\'ovil con mapa de instalaciones y reservas &
      Soluci\'on fiable y escalable; potencial referencia para el medio
      nacional \\
      \addlinespace
      Directorios municipales de espacios de recreaci\'on &
      Falta de visibilidad de los recintos &
      Plataforma centralizada con fichas de recintos y programaci\'on &
      Soluci\'on efectiva; validada por el p\'ublico y los gestores \\
      \addlinespace
      Redes sociales de escenarios locales, Bolivia &
      Difusi\'on fragmentada por recinto &
      Publicaci\'on individual en redes sociales &
      Soluci\'on parcial: no integra la oferta a nivel nacional \\
      \bottomrule
    \end{tabular}
  \end{tablaAPA}
\end{table}

\textbf{FUENTE:} Elaboraci\'on propia en base a la revisi\'on documental de
plataformas deportivas.

\subsection{\'Arbol de causas y efectos}

Para profundizar el an\'alisis de la situaci\'on problem\'atica se emplea la
t\'ecnica del \'arbol de problemas o \'arbol de causas y efectos. En este
diagrama, la parte central (tronco) representa el problema central; la parte
superior agrupa los efectos del problema y la parte inferior re\'une las
causas que lo originan \citep{arias2012}.

\begin{figure}[H]
  \centering
  \begin{tikzpicture}[
      node distance=1.2cm and 0.5cm,
      base/.style={
        rectangle, draw, thick, rounded corners=1.5mm, 
        minimum height=1.3cm, align=center, 
        font=\small, text width=4.5cm, inner sep=2mm
      },
      efecto/.style={base, fill=red!10, draw=red!70!black},
      problema/.style={base, fill=orange!15, draw=orange!80!black, text width=5.5cm, font=\small\bfseries},
      causa/.style={base, fill=blue!10, draw=blue!70!black},
      flecha/.style={-Stealth, thick, draw=black!60, line width=1.2pt}
    ]

    % 1. Tronco: problema central (Centro)
    \node[problema] (p) {INFORMACI\'ON SOBRE ESCENARIOS DEPORTIVOS DISPERSA Y SIN VISUALIZACI\'ON INTERACTIVA};

    % 2. Efectos del problema (Parte superior, ramificados)
    \node[efecto, above=1.5cm of p] (e2) {Subutilizaci\'on de la infraestructura deportiva existente};
    \node[efecto, left=0.5cm of e2] (e1) {Baja asistencia del p\'ublico a los eventos deportivos};
    \node[efecto, right=0.5cm of e2] (e3) {Dificultad para la promoci\'on del deporte y la actividad f\'isica};

    % 3. Causas del problema (Parte inferior, ramificadas)
    \node[causa, below=1.5cm of p] (c2) {Escasa digitalizaci\'on del sector deportivo};
    \node[causa, left=0.5cm of c2] (c1) {Ausencia de un registro centralizado de escenarios deportivos};
    \node[causa, right=0.5cm of c2] (c3) {Falta de herramientas de visualizaci\'on interactiva};

    % 4. Flechas (Causas -> Problema -> Efectos)
    
    % De las causas al problema
    \draw[flecha] (c1.north) -- (p.south);
    \draw[flecha] (c2.north) -- (p.south);
    \draw[flecha] (c3.north) -- (p.south);

    % Del problema a los efectos
    \draw[flecha] (p.north) -- (e1.south);
    \draw[flecha] (p.north) -- (e2.south);
    \draw[flecha] (p.north) -- (e3.south);

  \end{tikzpicture}
  \caption{\'Arbol de causas y efectos del problema identificado}
  \label{fig:arbol-causas-efectos}
\end{figure}

Las causas detectadas se han depurado agrupando aquellas que se encuentran
relacionadas entre s\'i y descartando las que se repet\'ian o que no afectaban
directamente a la poblaci\'on en la que se llevar\'a a cabo la investigaci\'on.
El an\'alisis causa--efecto expuesto permite pasar de la identificaci\'on del
problema a su formulaci\'on, que se presenta en la siguiente secci\'on.
